% TeX fragment extending the native227 section in the main report.
\section{A fourth saving: the final 226-state machine}
\label{sec:native226}

The selected public machine has 226 binary working states, with start 0,
dispatcher 82, and halt 226.  Its labels preserve ascending original K230 order
after deleting 65, 79, 140, and 175.  Its public table SHA-256 is
\begin{quote}\small\ttfamily
b58279b8b2e2cf029dd94e084d47cd6097772d7bce2a62653767895123d29599.
\end{quote}
The fully proved 227-state construction is retained as a checkpoint; the
226-state extension has its own source, literal tables, bijections and proofs.

A closed left reassociation at states $85/65$, applied at all ten occurrences
of the subtree, exposes an exact duplicate of the existing decoder at 98.
Both states have the rows
\[
 0\longmapsto 0R19,\qquad 1\longmapsto1R114.
\]
Consequently 65 can be deleted and all incoming targets 65 can be changed to
98.  Let $f(65)=98$ and $f(q)=q$ otherwise.  For every old control $q$ and read
symbol $b$, the checker verifies
\[
 \delta'(f(q),b)=(w,d,f(r))
 \quad\text{whenever}\quad \delta(q,b)=(w,d,r).
\]
All 454 old rows are included.  This finite equality is sufficient for a
one-step whole-tape simulation on arbitrary binary tapes and integer head
positions: the written bit and direction agree, and all other tape cells are
unchanged.  The map fixes start and halt and sends no working control to halt,
so deterministic induction yields halting equivalence with identical clocks
for the two tables in this quotient step.  No runtime guard or local tape
invariant is needed for the quotient.

The source DAG has 211 distinct nodes after this reassociation, compared with
212 for the 227-state source.  Leaf occurrences remain 5,121 and the guarded
source cases remain 6,656.  The initial reassociation creates four problematic
seven-zero address runs.  Five further closed reassociations repair them:
$101/209$ left, $101/208$ left, $226/224$ right, $205/134$ left, and $207/205$
right.  Their exact source boundaries and every literal row change are
recorded.  They preserve the native roles used by the earlier folds.  The final
source has maximum depth 41 and maximum zero run 6.

The unrepaired layout has a concrete instruction-protocol counterexample at
leaf 650, with address
\texttt{1110100100111000010000000}.  The zero branch of $CD(1)$, with all primary
registers zero, first spare zero, and gap zero, reaches a clean dispatcher after
658 steps with two wrong tape cells.  At relative positions 20 and 23 the
actual/expected values are $1/0$ and $0/1$.  The complete initial tape, exact
literal trace and both final tapes are retained.  The exact duplicate quotient
remains a valid theorem about the two literal intermediate tables; correctness
of the final encoding is established by the separate final protocol proof.

That final proof succeeds for all 6,656 guarded source frames, using 1,615,284
symbolic execution rules and all four finite-list induction schemas.  It
covers arbitrary natural primary register values and gap, arbitrary first
spare, and arbitrary finite spare tail.  All 452 rows are used.  Successor tape
and head position are exact, every block has positive duration, and the minimum
constant duration is 13.  Blank startup is checked literally in 409,835 steps,
with an additional dense canonical replay of the complete tape support.

Thus the 230-state input yields four state savings: the H loop fold, delayed
erasure, the reassignment of state 175's two decoder roles, and the newly
exposed duplicate decoder.  The source transformation and exact native proof
are composed with the frozen K230 consistency theorem.  The stated evidence
for the selected 226-state endpoint consists of ordinary mathematical proofs
and explicit induction rules; a fresh proof-assistant kernel check is a
separate claim and is not inferred from the endpoint's numerical state count.
